\chapter{분포 --- 데이터가 따르는 모양}
\label{ch:distributions}

\begin{keyidea}
분포를 가정하는 것은 \analogy{투영법}을 고르는 일이다. 편리해지지만 반드시
왜곡이 따르고, 어디가 왜곡되었는지는 검증해야만 알 수 있다.
\end{keyidea}

% =====================================================================
\section{정규분포}
\label{sec:normal}

\begin{inbrief}
  \item[Why] % TODO
  \item[When] % TODO
  \item[Where] % TODO
  \item[How] % TODO
\end{inbrief}

\begin{definition}[정규분포]
\label{def:normal}
\term{정규분포}(Normal Distribution) $X \sim \mathcal{N}(\mu, \sigma^2)$ 의
확률밀도함수(Probability Density Function)는
\[
  f(x) = \frac{1}{\sigma\sqrt{2\pi}}
         \exp\!\left(-\frac{(x-\mu)^2}{2\sigma^2}\right)
\]
\end{definition}

% TODO: 06_01_normal_fit.png, 06_02_normal_sigma_bands.png
% 68-95-99.7 규칙과, 이 데이터에서 실제로 몇 %가 들어갔는지(96.0%).

\begin{pitfall}[정규성 검정의 함정]
같은 표본에서 $\mu, \sigma$ 를 추정한 뒤 그 표본으로 KS 검정(Kolmogorov--Smirnov
Test)을 하면 p-value가 낙관적으로 나온다. 이 경우 Shapiro--Wilk 검정을 쓴다.
% TODO: calories 의 Shapiro p = 0.00038 -- 정규성이 기각된 사례로 정리.
\end{pitfall}

\section{지수분포}
\label{sec:exponential}
% TODO: $\lambda = 1/\text{mean}$, 06_04_exponential_fit.png

\section{균등분포}
\label{sec:uniform}
% TODO: 06_05_uniform_reference.png
% 주의: serving_size 는 균등분포가 아니다
%       (CDF 최대격차 0.153 > 정규분포의 0.101).

\section{베르누이와 이항분포}
\label{sec:binomial}
% TODO: 06_07_bernoulli_binomial.png

\section{누적분포함수로 비교하기}
\label{sec:cdf}
% TODO: 06_03_cdf.png -- 두 곡선의 최대 세로 격차가 KS 통계량 $D$.

\section{요약}
% TODO
